\section{What sampling hides}\label{sec:sampling}
The preceding tests attribute asymmetry to the sampled source. To relate that conclusion to thermodynamics, one must also account for motion between observations. The rotating model makes this remaining ambiguity explicit. Adding a full turn between samples changes the continuous motion without changing the observed transition.

Set $\nu=2\pi/\sqrt3$ in the model of Sec.~\ref{sec:main-cycle}. Rotation then completes a full turn per sample, and the transition is $I/4$. The sampled Gaussian source is reversible although its continuous-time probability current is nonzero. The detector-phase cancellation remains valid, but it cannot restore time order absent from the source samples.

For this reason, the entropy-production rate requires the continuous-time dynamics, not only their sampled law. For overdamped coordinates with even reversal parity, the stationary rate is~\cite{seifert2012}
\begin{equation}
 \sigma=\operatorname{tr}\{(A-D\Sigma^{-1})^{\mathsf T}
 D^{-1}(A-D\Sigma^{-1})\Sigma\}.
\end{equation}
It measures irreversibility in units of $k_B$ per unit time and vanishes at detailed balance for these models. The rotational rate is $6\nu^2/\lambda$. More generally,
\begin{equation}
 \nu_k=(2\pi/3+2\pi k)/\sqrt3,\qquad k\in\mathbb Z,
\end{equation}
gives the same sampled transition and innovation covariance for every integer $k$, while $6\nu_k^2/\lambda$ grows without bound. SM Sec.~S7.4 derives both the aliasing identity and the entropy rates. Thus even complete knowledge of the sampled law cannot determine continuous-time entropy production without additional dynamical or bandwidth information.
